
\documentclass[a4paper,12pt]{article}

\usepackage[utf8]{inputenc}
\usepackage[T1]{fontenc}
\usepackage{graphicx}
\usepackage{float}
\usepackage{geometry}
\usepackage{booktabs}
\usepackage{tikz}

\usetikzlibrary{arrows.meta,positioning}

\geometry{margin=2.5cm}

\setlength{\parindent}{0pt}
\setlength{\parskip}{0.5em}

\title{Bericht}
\author{Jiaqi Wang}
\date{\today}

\begin{document}

\maketitle


\section{Ausführen der Windows-Host-Software unter Ubuntu}

Die Host-Software wurde ursprünglich auf einem Windows-Computer mit Qt5 entwickelt.
Da die erzeugte Anwendung eine Windows-\texttt{.exe}-Datei ist, kann sie unter Ubuntu 22.04
nicht direkt nativ ausgeführt werden.

Zunächst wurde versucht, Qt5 direkt unter Ubuntu 22.04 zu installieren und die Host-Software
auf dem Linux-System erneut zu kompilieren. Dafür wurden die folgenden Befehle verwendet:

\begin{verbatim}
sudo apt-get install qt5-default
sudo apt update
sudo apt install build-essential qtbase5-dev qtbase5-dev-tools qt5-qmake
\end{verbatim}

Anschließend wurde überprüft, ob \texttt{qmake} korrekt installiert wurde:

\begin{verbatim}
qmake -v
\end{verbatim}

Danach wurde versucht, den Quellcode der Host-Software mit \texttt{cmake} neu zu kompilieren:

\begin{verbatim}
cd '/home/itmhiwi/jwang/itm-hiwi-main/host_pc/test'
mkdir build
cd build
cmake ..
\end{verbatim}

Dieser Ansatz ist jedoch aufwendiger, da die unter Windows entwickelte Qt5-Anwendung
möglicherweise zusätzliche Anpassungen für die Linux-Umgebung benötigt. Außerdem müssen
Abhängigkeiten, Projektkonfigurationen und plattformspezifische Einstellungen überprüft
und gegebenenfalls geändert werden.

Daher wurde anschließend eine einfachere und effizientere Methode gewählt: die Ausführung
der vorhandenen Windows-\texttt{.exe}-Datei unter Ubuntu mithilfe von Wine. Wine ist eine
Kompatibilitätsschicht, mit der viele Windows-Programme unter Linux ausgeführt werden können,
ohne dass eine vollständige Windows-Installation oder eine virtuelle Maschine benötigt wird.

\subsection{Installation von Wine}

Zunächst wird die 32-Bit-Architektur aktiviert, da viele Windows-Programme zusätzlich
32-Bit-Bibliotheken benötigen:

\begin{verbatim}
sudo dpkg --add-architecture i386
sudo apt update
\end{verbatim}

Anschließend werden Wine und die benötigten Zusatzpakete installiert:

\begin{verbatim}
sudo apt install wine wine64 wine32 winetricks
\end{verbatim}

Die installierte Wine-Version kann mit folgendem Befehl überprüft werden:

\begin{verbatim}
wine --version
\end{verbatim}

\subsection{Starten der Host-Software mit Wine}

Danach wird in das Verzeichnis gewechselt, in dem sich die Windows-\texttt{.exe}-Datei befindet:

\begin{verbatim}
cd '/home/itmhiwi/jwang/MotroController_APP_V1'
cd pkg/
\end{verbatim}

Die Host-Software kann anschließend mit Wine gestartet werden:

\begin{verbatim}
wine MotroController.exe
\end{verbatim}


\end{document}
